\chapter{Chapter heading}
\label{chap:one}
% labels don't require specific prefixes, but I find it conceptually helpful
% and of course the labels themselves are optional unless actually referred to

\lettrine{Y}{ou can insert}
dropped capitals, also called \emph{drop caps}, with the \texttt{lettrine} package.\falsepar The \verb|\falsepar| macro can be used as demonstrated if the first paragraph is very short and you need to start a new paragraph before passing the bottom of the dropcap. Search for the package name in \texttt{uwthesis.sty} file if you want to change the font, and see the package documentation for more many options (\url{https://ctan.org/pkg/lettrine}).



\section{First level section heading}
\label{sec:one}

\subsection{Second level (sub)section heading}
\label{sec:two}

Section headings should almost certainly be sentence case (\url{https://practicaltypography.com/headings.html}). Whether Chapter headings/titles should use sentence case or Title Case is much more open to personal preference, various style guides suggest differently.


\subsubsection{Third level (subsub)section heading}
\label{sec:three}

There is a third level of headings (and even lower level \verb|\paragraph| and \verb|\subparagraph| commands, see the \texttt{memoir} class manual) but you may want to consider whether such deep nesting is really required.


\section{Cleveref}
\label{sec:cleveref}

The \texttt{cleveref} package is particularly helpful for referring to items in your manuscript such as figures/tables/sections/chapters/etc. You generally use the \verb|\cref| macro for most references, with the capitalised \verb|\Cref| used at the beginning of a sentence. The benefit is that the package automagically determines the correct type(s) of the object(s) and can handle references to multiple items at once. It also creates hyperlinks to the referenced object in conjunction with the \texttt{hyperref} package. To give a couple quick examples:

\Cref{chap:one} is the first chapter containing \cref{sec:one,sec:two} and \cref{fig:example}.

As usual, check the package documentation (\url{https://ctan.org/pkg/cleveref}) for more options, particularly if you want to change the form used for different names (e.g.~Fig.\ vs figure).


\section{Abbreviations}
\label{sec:abbreviations}

To use abbreviations in the main text, use the un-starred \texttt{acro} package macros for all your \acp{TLA} and \acp{DoF}, and don't forget your \ac{SCUBA}. Using this package has the advantage that it will properly expand the abbreviation upon its first (un-starred) usage, so you don't have to manually make changes if you add or remove uses. With the \texttt{hyperref} package, the abbreviation will also serve as a link back to the corresponding entry in the list of abbreviations.

\begin{figure}
    \includegraphics[width=0.9\textwidth]{example-image}
    \caption[Short caption to go in list of figures]%
    {Longer descriptive caption to go underneath the figure. If you use abbreviations in captions (particularly in the short captions) then use the starred \texttt{acro} package macros, and probably the ``short'' versions to ensure it is abbreviated, e.g.~\acs*{TLA}.}
    \label{fig:example}
\end{figure}



\section{Tables}

Tables can/should be made using the \texttt{booktabs} package macros \verb|\toprule|, \verb|\midrule|, and \verb|\bottomrule|, as shown in the examples. This package is included by the \texttt{memoir} class we are using, so it does not need to be included separately. See the package documentation for more details and typographical hints on tables (\url{https://ctan.org/pkg/booktabs}). In particular, avoid the use of vertical rules (i.e.~no vertical lines), numerical columns should generally right aligned (especially to align decimal points), while text columns are generally left aligned. Centre multi-column headers.

For a very brief overview of some tips and tricks, have a gander at Markus Püschel's \emph{Small Guide to Making Nice Tables}: \url{ https://people.inf.ethz.ch/markusp/teaching/guides/guide-tables.pdf}.

\begin{table}
    \centering
% the @{} removes the white space from the edge of the table
    \begin{tabular}{@{}llr@{}}
        \toprule
        Animals & Colours & Numbers \\
        \midrule
        dog & blue & 3.14 \\
        cat & green & 2.71 \\
        \bottomrule
    \end{tabular}
    \caption{A simple table example}
    \label{tab:my_label}
\end{table}

\begin{table}
    \centering
% Overleaf may parse an error here, but it's fine; the >{$}c<{$} inserts dollar
% signs before and after the contents of the 1st column to set it in math mode.
% Very useful if you have a full column of math, rather than surrounding every
% individual entry with $$
    \begin{tabular}{@{}>{$c<$}rrcrr@{}}
        \toprule
        \multirow{3}{*}[-4pt]{\begin{tabular}{@{}c@{}}Multi-row\\heading\end{tabular}} & \multicolumn{5}{c}{Multi-column super heading} \\
        \cmidrule(l){2-6}
        & \multicolumn{2}{c}{Cats} && \multicolumn{2}{c}{Dogs} \\
        \cmidrule(lr){2-3} \cmidrule(l){5-6}
        & count & \% && count & \% \\
        \midrule
        1 & 6646 & 57.7 && 53887 & 5.1 \\
        (1,4] & 1897 & 16.5 && 49152 & 4.7 \\
        \bottomrule
    \end{tabular}
    \caption{A more complex table example}
    \label{tab:complex}
\end{table}
